\documentclass[11pt]{article}

\usepackage[T1]{fontenc}
\usepackage{mlmodern}
\usepackage[margin=1in]{geometry}
\usepackage{amsmath,amssymb,amsthm}
\usepackage{booktabs}
\usepackage{array}
\usepackage{xurl}
\usepackage{hyperref}

\hypersetup{
  colorlinks=true,
  linkcolor=blue,
  citecolor=blue,
  urlcolor=blue,
  pdftitle={An Exact Arithmetic Profile of X0(137): Fricke Decomposition, Newform Orbits, Period Maps, and Dedekind Data},
  pdfauthor={Luqman Mahmood},
  pdfsubject={Modular curves, Fricke involutions, newforms, period maps, class groups, and Dedekind sums},
  pdfkeywords={X0(137), Fricke involution, modular symbols, newforms, period maps, class groups, Dedekind sums}
}
\ifdefined\pdftrailerid
  \pdftrailerid{0123456789ABCDEF0123456789ABCDEF}
\fi

\newtheorem{theorem}{Theorem}[section]
\newtheorem{proposition}[theorem]{Proposition}
\newtheorem{corollary}[theorem]{Corollary}
\theoremstyle{definition}
\newtheorem{definition}[theorem]{Definition}
\theoremstyle{remark}
\newtheorem{remark}[theorem]{Remark}

\newcommand{\C}{\mathbf{C}}
\newcommand{\Q}{\mathbf{Q}}
\newcommand{\Z}{\mathbf{Z}}
\DeclareMathOperator{\Hom}{Hom}
\DeclareMathOperator{\Cl}{Cl}
\DeclareMathOperator{\tr}{tr}
\DeclareMathOperator{\ord}{ord}

\title{An Exact Arithmetic Profile of $X_0(137)$\\
  \large Fricke Decomposition, Newform Orbits, Period Maps, and Dedekind Data}
\author{Luqman Mahmood\thanks{Copyright 2026 Luqman Mahmood. This paper is licensed under Creative Commons Attribution 4.0 International (CC BY 4.0); see \url{https://creativecommons.org/licenses/by/4.0/}. Software is excluded from this grant.}}
\date{30 September 2026}

\begin{document}

\maketitle

\begin{abstract}
We determine the exact arithmetic profile of the modular curve $X_0(137)$.
The curve has genus $11$; the Fricke involution acts on its weight-two cusp
space with eigenspace dimensions $(4,7)$ and trace $-3$, and the quotient
curve has genus $4$.  The class group of discriminant $-548$ is cyclic of
order $8$, generated by the form $(3,2,46)$, whose full power cycle we
display.  The weight-two newspace consists of two Galois orbits of degrees
$4$ and $7$.  We give the minimal polynomials of the Hecke eigenvalue $a_2$
for both orbits and distinguish them from the coefficient-field defining
polynomials recorded in the LMFDB.  An exact rational period-map computation
shows that the winding element has nonzero image for the degree-seven orbit,
proving analytic rank zero for each of its conjugates.  For the degree-four
orbit the central value vanishes exactly; numerical evaluation of the first
derivative suggests rank one but is not used as a theorem.

We also study the subspaces $V^\pm$ of cusp forms vanishing at both cusps,
derive their cross-Hom dimension, and use its value $18$ at level $137$
as a worked inverse problem.  Finally, we compute the classical
Dedekind sums $s(h,137)$, determine their zero set $\{37,100\}$, and reduce
five ratio classifications to residue-class quadratics using Dedekind
reciprocity.  All theorem-level computations use exact arithmetic; the
underlying finite data are displayed in the text or described in the
appendix.
\end{abstract}

\noindent\textbf{2020 Mathematics Subject Classification.}
Primary 11F30, 11F20; Secondary 11F67, 11R29.

\noindent\textbf{Keywords.}
Modular curves, Fricke involutions, modular symbols, newforms, period maps,
class groups, Dedekind sums.

\section{Introduction}

The modular curve $X_0(137)$ provides a compact setting in which several
classical arithmetic structures can be determined exactly and compared in a
single normalization.  We study its genus and Fricke decomposition, the
proper class group of discriminant $-548$, the Galois orbits in its
weight-two newspace, and the classical Dedekind sums at modulus $137$.  Two
points require particular care.  First, a polynomial defining a Hecke field
need not be the minimal polynomial of the specified coefficient $a_2$.
Second, numerical evidence for a derivative of an $L$-function is not an
exact rank certificate.  Keeping these distinctions explicit yields a sharp
arithmetic profile without enlarging the theorem set beyond the available
proofs.

The Fricke decomposition also gives two naturally different Hom spaces.  The
full eigenspaces and the cusp-vanishing hyperplanes satisfy
\[
  \dim\Hom(S_p^+,S_p^-)=d_+(p)d_-(p),\qquad
  \dim\Hom(V^+,V^-)=(d_+(p)-1)(d_-(p)-1).
\]
We prove the second formula for prime level and determine the primes for which
its value equals the target $18$ read from level $137$.
This is a worked inverse problem inside the stated modular-curve family,
not an independent prediction of the target or a preferred level.

The Dedekind-sum results form a second exact classification.  Reciprocity
turns each equation $s(h_1,p)/s(h_2,p)=a/b$ into finitely many linear or
quadratic equations indexed by residue classes.  This proves five fixed-pair
prime loci and, after an exhaustive finite calculation, the simultaneous
$h\leq28$ locus.  The pairs, ratios, and cutoff are hypotheses of these
statements rather than invariants selected by the theory.  In particular,
the target fractions are taken from level $137$ before asking whether other
prime levels reproduce them.  This illustrates an exact decision procedure;
it supplies no independent explanation of the selected inputs.
No physical interpretation is attached to these results.

All theorem-level computations are over exact rings or exact rational
arithmetic.  The larger matrix and residue-class calculations have
machine-readable data and independent verification routines described in
Appendix~A.  The only numerical statement retained in the discussion is the
suggested rank-one behavior of the degree-four newform orbit, which remains
explicitly conditional.

Section~2 fixes conventions and the genus bound.  Sections~3--7 treat the
Fricke, class-group, newform, cusp-vanishing, and fixed-target Hom
calculations.  Sections~8 and~9 give the Dedekind values and prime-locus
classifications, followed by concluding questions and data availability.

\section{Conventions and the prime-level genus bound}

Let $p$ be prime and write
\[
  X_p=X_0(p), \qquad S_p=S_2(\Gamma_0(p)).
\]
For $p>3$, the curve has two cusps, denoted $0$ and $\infty$.  Let $w_p$
be the Fricke involution on $X_p$, and let $W_p=w_p^*$ be its pullback on
\[
  S_p \simeq H^0(X_p,K_{X_p});
\]
the weight-two differential identification is standard
\cite[Chapter~3, especially p.~36]{Stein2007}.  We use
\[
  S_p^{\varepsilon}=\ker(W_p-\varepsilon), \qquad
  d_\varepsilon(p)=\dim_{\C} S_p^\varepsilon,
  \qquad \varepsilon\in\{+1,-1\}.
\]
On the upper half-plane this convention is
\begin{equation}\label{eq:fricke-pullback}
  (W_p f)(z)=\frac{1}{p z^2}f\!\left(-\frac{1}{p z}\right),
  \qquad
  w_p^*(f(z)\,dz)=(W_p f)(z)\,dz.
\end{equation}
The modular-symbol star sign is distinct from the Fricke eigenvalue
$W_p=\pm1$ used below.  We determine $d_+$ and $d_-$ as the nullities of
$W_p-I$ and $W_p+I$ on the cuspidal modular-symbol realization, so the signs
in the tables are fixed by the operator itself rather than by an interface
label.

For prime $p>3$, the exact genus formula used in the finite-locus proof is
\begin{equation}\label{eq:genus}
  g(X_p)=\frac{p+1}{12}-\frac{e_2(p)}{4}-\frac{e_3(p)}{3},
  \qquad
  e_2(p)=1+\left(\frac{-1}{p}\right),\quad
  e_3(p)=1+\left(\frac{-3}{p}\right).
\end{equation}
This is the prime specialization of the general $X_0(N)$ formula recorded by
Ogg \cite[p.~455]{Ogg1974}.  Since $e_2,e_3\in\{0,2\}$,
\begin{equation}\label{eq:genusbound}
  g(X_p)\geq \frac{p+1}{12}-\frac{1}{2}-\frac{2}{3}
  =\frac{p-13}{12}.
\end{equation}

For later transparency, write $n_p$ for the number of fixed points of $w_p$.
Riemann--Hurwitz and the identification of invariant differentials with
differentials on the quotient give
\begin{equation}\label{eq:fixed-trace}
  \tr(W_p)=1-\frac{n_p}{2},
  \qquad
  d_+(p)=\frac{g(X_p)+\tr(W_p)}{2},
  \qquad
  d_-(p)=\frac{g(X_p)-\tr(W_p)}{2}.
\end{equation}
Indeed, if $Y=X_p/\langle w_p\rangle$, then
$2g(X_p)-2=2(2g(Y)-2)+n_p$, while $g(Y)=d_+(p)$.

\section{Genus and Fricke decompositions}

\begin{theorem}[Genus and weight-two dimension at level $137$]\label{thm:genus-137}
For $X=X_0(137)$,
\[
  g(X)=11,
  \qquad
  \dim_{\C}S_2(\Gamma_0(137))=11.
\]
\end{theorem}

\begin{proof}
In \eqref{eq:genus}, $e_2(137)=2$ and $e_3(137)=0$, giving $g=11$.
The weight-two differential identification gives
$\dim_{\C}S_2(\Gamma_0(137))=g=11$.
\end{proof}

\begin{theorem}[Fricke data at level $137$]\label{thm:fricke-137}
For $X=X_0(137)$, the Fricke eigenspace dimensions, trace, and quotient genus
are
\[
  (d_+,d_-)=(4,7), \qquad \tr(W_{137})=-3,
  \qquad g\bigl(X_0(137)/\langle w_{137}\rangle\bigr)=4.
\]
\end{theorem}

\begin{proof}
Ogg's fixed-point/class-number formula gives $n_{137}=8$
\cite[p.~454]{Ogg1974}.  Equations \eqref{eq:fixed-trace} and
Theorem~\ref{thm:genus-137} give $\tr(W_{137})=-3$ and
$(d_+,d_-)=(4,7)$; Riemann--Hurwitz gives quotient genus $4$
\cite{StacksRiemannHurwitz}.

For completeness, invariant differentials really do descend through the
branched quotient.  Away from the fixed points this is ordinary etale descent.
At a fixed point choose a local coordinate $z$ with $w_{137}(z)=-z$.  If
$f(z)\,dz$ is invariant, then $f(z)=-f(-z)$, so
$f(z)=z g(z^2)$.  With $t=z^2$,
\[
  f(z)\,dz=\frac{1}{2}g(t)\,dt,
\]
which is the pullback of a holomorphic differential in the quotient
coordinate.  Hence the invariant differential space has dimension equal to
the quotient genus.

Independently, exact modular-symbol calculation gives an involution with
characteristic polynomial
\[
  \det(TI-W_{137})=(T-1)^4(T+1)^7.
\]
This reproduces the split and trace directly in the convention
\eqref{eq:fricke-pullback}.
\end{proof}

\begin{proposition}[Neighboring genus-$11$ rows]\label{prop:neighbor-fricke-splits}
The three primes $131$, $137$, and $139$ all give genus $11$, with
\[
\begin{array}{c|ccc}
  p & d_+(p) & d_-(p) & \tr(W_p)\\
  \hline
  131 & 1 & 10 & -9\\
  137 & 4 & 7  & -3\\
  139 & 3 & 8  & -5
\end{array}
\]
In particular, the row at $139$ is $(3,8)$, not $(4,7)$.
\end{proposition}

\begin{proof}
Ogg's formula gives
\[
  (n_{131},n_{137},n_{139})=(20,8,12).
\]
Equation \eqref{eq:fixed-trace} gives traces
$(-9,-3,-5)$ and hence the displayed splits.  Exact modular-symbol
calculation over the coefficient ring gives the nullities of $W_p-I$ and
$W_p+I$.  In each row $W_p^2=I$, $d_++d_-=g(X_p)$, and the genus agrees with
\eqref{eq:genus}.  The displayed trace is then $d_+-d_-$.
\end{proof}

\begin{corollary}[Full cross-Fricke Hom]\label{cor:full-hom-28}
At level $137$,
\[
  \dim_{\C}\Hom(S_{137}^+,S_{137}^-)=4\cdot 7=28.
\]
\end{corollary}

\begin{proof}
For finite-dimensional complex vector spaces $A$ and $B$,
$\dim\Hom(A,B)=(\dim A)(\dim B)$.  Apply Theorem~\ref{thm:fricke-137}.
\end{proof}

\section{The class group of discriminant \texorpdfstring{$-548$}{-548}}

We write $\Cl(-548)$ for the proper class group of primitive positive-definite
binary quadratic forms of discriminant $-548$.  This is equivalently the ideal
class group of $\Q(\sqrt{-137})$, whose fundamental discriminant is $-548$.

\begin{theorem}\label{thm:class-group-minus-548}
The group $\Cl(-548)$ is cyclic of order $8$.  The reduced primitive form
\[
  Q=(3,2,46)
\]
has exact order $8$ and generates the group.  Under the form--ideal
correspondence, it is the class of $(3,w+2)$ in
$\Q(w)$, where $w^2+137=0$.
\end{theorem}

\begin{proof}
Exact reduction and composition of forms gives the following power cycle.
\[
\begin{array}{c|c}
  n & \text{reduced representative of }Q^n\\
  \hline
  1&(3,2,46)\\
  2&(9,8,17)\\
  3&(6,-2,23)\\
  4&(2,2,69)\\
  5&(6,2,23)\\
  6&(9,-8,17)\\
  7&(3,-2,46)\\
  8&(1,0,137)
\end{array}
\]
Every row has discriminant $-548$.  The first seven rows are nonprincipal,
the eighth is principal, and the eight rows are precisely the complete list of
reduced primitive classes.  Thus $Q$ has exact order $8$, the class number is
$8$, and $Q$ generates.  The displayed cycle is already an explicit
check; independent ideal-class and quadratic-form computations give the
same invariant factor $[8]$.  The form and class-group algorithms are documented in
\cite{SageQuadraticForms}.

For the corresponding ideal $I=(3,w+2)$, exact ideal arithmetic gives
\[
  I^8=(56-5w), \qquad N(I^8)=3^8=6561.
\]
The opposite standard choice in the form--ideal correspondence replaces this
generator by its inverse; its order and the cyclic group statement are
unchanged.
\end{proof}

\section{Level-137 newforms and the coefficient-polynomial distinction}

At prime level $137$ and weight $2$, the oldspace inherited from level $1$ is
zero, so the cusp space is new.  For a normalized Fricke eigenform write
$W_{137}f=\lambda_{137}(f)f$.  We use the root-number convention
\[
  \epsilon(f)=-\lambda_{137}(f)
\]
for normalized trivial-character newforms of this prime level.  Newform orbit
and Atkin--Lehner conventions are reviewed in
\cite[Chapter~9]{Stein2007} and
\cite[Sections~4.2 (p.~7), 4.5 (pp.~9--11), and 8.3 (pp.~28--29)]{BestEtAl2021}.
The $W_{137}$ column below records the scalar $\lambda_{137}(f)$, while
$\epsilon$ is the root number under the displayed sign convention; neither
column is the modular-symbol star sign.

For an orbit with conjugates $f^\sigma$, ``analytic rank per conjugate'' means
$\ord_{s=1}L(f^\sigma,s)$.  The order of the orbit-product $L$-function is the
sum over conjugates.  We state both when a distinction matters.

\begin{theorem}\label{thm:newform-orbits}
The level-$137$, weight-two, trivial-character newspace has exactly two Galois
newform orbits.  Their degrees, Fricke signs, and root numbers are
\[
\begin{array}{c|cc}
  \text{orbit degree} & W_{137} & \epsilon\\
  \hline
  4 & +1 & -1\\
  7 & -1 & +1.
\end{array}
\]
The actual minimal polynomials of $a_2$ are, respectively,
\begin{align}
  m_{4,a_2}(x)&=x^4+3x^3-4x-1,\label{eq:a2four}\\
  m_{7,a_2}(x)&=x^7-10x^5+28x^3+3x^2-19x-7.\label{eq:a2seven}
\end{align}
The degree-seven orbit has analytic rank $0$, certified by algebraic central
nonvanishing.  Exact central vanishing holds for the degree-four orbit.
\end{theorem}

\begin{proof}
An exact $11\times11$ matrix for $T_2$ on the sign-$+1$
cuspidal modular-symbol realization has characteristic polynomial
\[
  \det(xI-T_2\mid S_2(\Gamma_0(137)))
  =m_{4,a_2}(x)m_{7,a_2}(x).
\]
The calculation gives an integral cyclic change-of-basis matrix $P$, with
$\det P=16$, and the exact identity $T_2P=PC$ for the companion matrix $C$ of
this product.  The two factors are irreducible over $\Q$, as witnessed by
irreducible reductions modulo $2$ and $3$, respectively.  Their degrees sum
to $11$, and their exact Atkin--Lehner eigenvalues reproduce the $(4,7)$
Fricke split.

Let $f$ be a normalized representative of the degree-seven orbit, put
$A_7=A_f$, and let $e=\{0,\infty\}$ be its winding symbol.  For every
conjugate $f^\sigma$, the integration pairing satisfies
\begin{equation}
  L(f^\sigma,1)=-2\pi i\langle f^\sigma,e\rangle
        =-2\pi i\int_0^{i\infty} f^\sigma(z)\,dz.
  \label{eq:winding-central-value}
\end{equation}
This is the weight-two case of the modular-symbol special-value formula
\cite[Equation~(8.5.5)]{Stein2007}.  A rational period mapping has the same
kernel as the complex period mapping
\cite[Section~10.6]{Stein2007}.  SageMath documents this kernel property and
the resulting exact central-vanishing test for simple factors
\cite{SageModularSymbolsPeriods,SageAbvarLseries}.
For a simple modular abelian variety $A_f$, the resulting exact criterion is
\begin{equation}
  L(A_f,1)=0
  \quad\Longleftrightarrow\quad
  \Theta_f(e)=0,
  \label{eq:winding-vanishing-criterion}
\end{equation}
where $\Theta_f$ is any rational period mapping for the newform quotient.

In the stated rational bases, the exact $23\times14$ rational period map for
the simple modular abelian variety $A_7$ sends the winding symbol $e$ to
\[
  (-1,0,0,0,0,0,0,0,0,0,0,0,0,0).
\]
The vector is nonzero, so
\eqref{eq:winding-vanishing-criterion} proves $L(A_7,1)\neq0$.  Since
$L(A_7,s)=\prod_\sigma L(f^\sigma,s)$, every conjugate central value is
nonzero; each conjugate and the orbit product therefore have analytic order
zero at $s=1$.  The persisted rational matrix and basis data allow the
nonzero period-map image to be checked directly by exact matrix
multiplication.

For the degree-four orbit, the algebraic route gives exact central vanishing
for every conjugate, consistently with root number $-1$.
\end{proof}

\begin{remark}
High-precision evaluations at $96$ and $160$ bits give a stable nonzero first
derivative at every real embedding, suggesting rank one per conjugate and
orbit-product order $4$.  Because these evaluations carry no certified error
interval, both order-one statements remain conditional.  The distinction
between numerical rank calculations and rigorous certification is discussed in
\cite[Section~9.4, pp.~36--38]{BestEtAl2021}.
\end{remark}

\subsection{Coefficient-field polynomials are not \texorpdfstring{$a_2$}{a2} polynomials}

The following two polynomials occur in the LMFDB records for the two
orbits \cite{LMFDB137a,LMFDB137b}:
\begin{align}
  F_4(x)&=x^4-x^3-3x^2+x+1,\label{eq:fieldfour}\\
  F_7(x)&=x^7-10x^5+28x^3-3x^2-19x+7.\label{eq:fieldseven}
\end{align}
They are valid defining polynomials for chosen presentations of the two Hecke
fields.  They are \emph{not} the minimal polynomials of the specified
coefficient $a_2$.  The latter are \eqref{eq:a2four} and
\eqref{eq:a2seven}.  A field presentation and the minimal polynomial of a
particular element of that field are different data, even when that element
happens to generate the field.  This resolves the apparent polynomial
discrepancy; no coefficient-level conclusion is inferred from either
presentation.  Explicitly, $F_7(x)=-m_{7,a_2}(-x)$, so the displayed
degree-seven field generator is the negative of the specified Hecke generator.
Thus $F_4,F_7$ are coefficient-field defining polynomials recorded by the
database, whereas $m_{4,a_2},m_{7,a_2}$ are the minimal polynomials of the
specified Hecke coefficient $a_2$.  The two kinds of polynomial must not be
substituted for one another in an arithmetic statement.

\section{The canonical cusp-vanishing construction}

Let $p>3$ be prime, let $X=X_p$, and set
\[
  D=(0)+(\infty), \qquad V=H^0(X,K_X-D)\subset H^0(X,K_X)=S_p.
\]
Thus $V$ consists of holomorphic differentials having an additional zero at
both cusps.  Since $w_p$ exchanges $0$ and $\infty$, $V$ is Fricke-stable; put
$V^\varepsilon=V\cap S_p^\varepsilon$.

If
\[
  f(q)=\sum_{n\geq 1}a_n(f)q^n\in S_p,
\]
then the differential associated to $f$ is
\[
  f(q)\,\frac{dq}{2\pi i q}
  =\frac{1}{2\pi i}\bigl(a_1(f)+a_2(f)q+\cdots\bigr)dq.
\]
It vanishes at $\infty$ exactly when $a_1(f)=0$.

\begin{theorem}[Cusp-vanishing Fricke blocks]\label{thm:cusp-reduced-hom}
For every prime $p>3$, the space $V$ is Fricke-stable.  If
$d_\varepsilon(p)>0$, then
\[
  V^\varepsilon
  =\ker\bigl(a_1:S_p^\varepsilon\longrightarrow\C\bigr),
  \qquad
  \dim V^\varepsilon=d_\varepsilon(p)-1.
\]
If both Fricke blocks are nonzero, then
\begin{equation}\label{eq:reducedhom}
  \dim\Hom(V^+,V^-)
  =(d_+(p)-1)(d_-(p)-1).
\end{equation}
At $p=137$,
\[
  (\dim V^+,\dim V^-)=(3,6),
  \qquad \dim\Hom(V^+,V^-)=18.
\]
If either Fricke block is zero, the corresponding $V^\varepsilon$ and the
cross-Hom space are zero.
\end{theorem}

\begin{proof}
Because $w_p$ exchanges the two points of $D$, it preserves the line bundle
$K_X-D$ and hence $V$.  On a $W_p$-eigenvector, vanishing at one cusp is
equivalent to vanishing at the other.  The local $q$-expansion above therefore
gives
\[
  V^\varepsilon=\ker(a_1|S_p^\varepsilon).
\]

It remains to show that $a_1$ is not the zero functional on a nonzero Fricke
block.  The prime-level weight-two cusp space is new because
$S_2(\mathrm{SL}_2(\Z))=0$.  The newform decomposition is compatible with
Fricke, and every nonzero constituent contains a normalized eigenform with
$a_1=1$; equivalently, the perfect Hecke pairing and
$a_1(T_n f)=a_n(f)$ give the same conclusion
\cite[Lemma~3.22 and Proposition~3.24]{Stein2007}.  Thus the kernel has
codimension one whenever the block is nonzero.  Equation
\eqref{eq:reducedhom} is finite-dimensional linear algebra.  Substituting
$(d_+,d_-)=(4,7)$ gives $(3,6)$ and $18$.  If a block is zero, its intersection
with $V$ is zero and the cross-Hom space is zero; the bare signed product in
\eqref{eq:reducedhom} is not used in that edge case.
\end{proof}

\begin{remark}[Hecke boundary]
The space $V$ is Fricke-stable, but when $V\neq0$ it is not stable under every
Hecke operator.  Indeed, if every $T_n$ preserved $V$, then for $f\in V$ one
would have $a_n(f)=a_1(T_n f)=0$ for every $n$, so the $q$-expansion principle
would force $f=0$.  Thus the construction uses the canonical hyperplanes
$\ker(a_1)$ as vector spaces.  No quotient Hecke algebra, Hecke-stable reduced
block, or canonical complementary line is asserted.
\end{remark}

\begin{remark}[Full and cusp-vanishing Hom spaces are distinct]\label{rem:distinct-hom}
At $p=137$,
\[
  \dim\Hom(S_{137}^+,S_{137}^-)=28,
  \qquad
  \dim\Hom(V^+,V^-)=18.
\]
The second space is defined from canonical hyperplanes inside the two Fricke
blocks.  There is a canonically defined restriction sequence
\[
  0\longrightarrow
  \Hom(S_{137}^+/V^+,S_{137}^-)
  \longrightarrow E
  \longrightarrow \Hom(V^+,V^-)
  \longrightarrow0,
\]
where
\[
  E=\{T\in\Hom(S_{137}^+,S_{137}^-):T(V^+)\subseteq V^-\}.
\]
The last arrow restricts a map to $V^+$.  It is surjective because linear maps
extend after a choice of complements, but no such choice is canonical.  Hence
the exact sequence has no selected canonical splitting: there is no asserted
canonical inclusion of the $18$-dimensional Hom space into the
$28$-dimensional one, and the two spaces are not identified.
\end{remark}

\section{A worked cross-Fricke inverse problem}

We now fix the modular-curve family, weight, character, cusp divisor, Fricke
convention, and numerical target, and ask which prime levels attain that
target.  Here that target is the already computed value at $p=137$.

\begin{theorem}\label{thm:target-conditioned-locus}
Fix the family $X_0(p)$ for prime $p$, weight $2$, trivial character, the
Fricke convention above, the divisor $D=(0)+(\infty)$ for $p>3$, and the
target
\[
  \dim\Hom(V^+,V^-)=18.
\]
Then $p=137$ is the unique prime attaining the target.  For $p=2$ and $p=3$,
the weight-two cusp space is zero and the reduced cross-Hom dimension is set
to zero.
\end{theorem}

\begin{proof}
First suppose $p>3$ and both Fricke blocks are nonzero.  Put
$u=d_+(p)-1$ and $v=d_-(p)-1$.  If $uv=18$, the positive factor pairs up to
interchange are
\[
  (u,v)\in\{(1,18),(2,9),(3,6)\}.
\]
Because $g(X_p)=d_+(p)+d_-(p)=u+v+2$, the only possible genera are
\[
  21,\quad 13,\quad 11.
\]
The bound \eqref{eq:genusbound} then forces $p\leq265$.  Exact enumeration of
the primes in this bounded range gives the complete list of rows having one of
the three forced genera shown below.  The genus and Fricke nullities are
computed over exact coefficient rings.
The fixed-point counts are the corresponding exact Ogg values, and
\eqref{eq:fixed-trace} independently turns them into the same traces and
splits.
\[
\begin{array}{c|c|c|c|c|c}
  p & g(X_p) & n_p & \tr(W_p) & (d_+,d_-) & \dim\Hom(V^+,V^-)\\
  \hline
  131 & 11 & 20 & -9  & (1,10) & 0\\
  137 & 11 & 8  & -3  & (4,7)  & 18\\
  139 & 11 & 12 & -5  & (3,8)  & 14\\
  163 & 13 & 4  & -1  & (6,7)  & 30\\
  251 & 21 & 28 & -13 & (4,17) & 48\\
  257 & 21 & 16 & -7  & (7,14) & 78
\end{array}
\]
The only target row is $137$.  If one Fricke block is zero, the cross-Hom
dimension is zero by Theorem~\ref{thm:cusp-reduced-hom}.  Finally,
$g(X_0(2))=g(X_0(3))=0$, so their weight-two cusp spaces and cross-Hom spaces
are zero.  This exhausts all primes.
\end{proof}

\begin{remark}
The uniqueness is relative to the stated family and the fixed target $18$;
the theorem does not select that target independently.
\end{remark}

\section{Classical Dedekind sums at modulus 137}\label{sec:classical-dedekind}

\begin{definition}[Classical Dedekind sum]\label{def:classical-dedekind-sum}
For an integer $h$ and a positive integer $k$, define
\[
  s(h,k)=\sum_{j=1}^{k-1}
  \left(\!\left(\frac{j}{k}\right)\!\right)
  \left(\!\left(\frac{hj}{k}\right)\!\right),
\]
where
\[
  ((x))=
  \begin{cases}
    x-\lfloor x\rfloor-\frac12,&x\notin\Z,\\
    0,&x\in\Z.
  \end{cases}
\]
The finite sum is defined without a coprimality hypothesis.  Coprimality enters
the inversion and reciprocity identities below, not the definition.
\end{definition}

\begin{proposition}[Classical identities]\label{prop:dedekind-identities}
For integers $h$ and positive integers $k$:
\begin{enumerate}
  \item $s(h+k,k)=s(h,k)$;
  \item $s(-h,k)=-s(h,k)$;
  \item if $\gcd(h,k)=1$ and $hh'\equiv1\pmod{k}$, then
        $s(h',k)=s(h,k)$; and
  \item if $h,k>0$ and $\gcd(h,k)=1$, then
  \begin{equation}\label{eq:dedekind-reciprocity}
    s(h,k)+s(k,h)
    =-\frac14+\frac1{12}\left(\frac hk+\frac kh+\frac1{hk}\right).
  \end{equation}
\end{enumerate}
\end{proposition}

\begin{proof}
The first two assertions follow directly from the sawtooth function, and the
third follows by permuting the nonzero residue classes modulo $k$.  The fourth
is the classical Dedekind reciprocity law; see
\cite[Chapter~2]{RademacherGrosswald1972}.
\end{proof}

\begin{proposition}[Selected exact values at modulus 137]
\label{cert:dedekind-values-137}
The following values and reduced-numerator factorizations are exact.
\[
\begin{array}{c|c|l@{\qquad}c|c|l}
h&s(h,137)&\text{numerator}&h&s(h,137)&\text{numerator}\\
\hline
3  &495/137&3^2\cdot5\cdot11 &13&45/137 &3^2\cdot5\\
7  &180/137&2^2\cdot3^2\cdot5&14&108/137&2^2\cdot3^3\\
8  &102/137&2\cdot3\cdot17   &15&18/137 &2\cdot3^2\\
9  &120/137&2^3\cdot3\cdot5 &18&44/137 &2^2\cdot11\\
11 &140/137&2^2\cdot5\cdot7 &19&10/137 &2\cdot5\\
22 &20/137 &2^2\cdot5        &28&90/137 &2\cdot3^2\cdot5\\
37 &0      &0                 &100&0     &0
\end{array}
\]
In particular,
\begin{align}
  \frac{s(22,137)}{s(7,137)}&=\frac19,&
  \frac{s(19,137)}{s(9,137)}&=\frac1{12},&
  \frac{s(13,137)}{s(7,137)}&=\frac14,\label{eq:dedekind-five-ratios-a}\\
  \frac{s(14,137)}{s(28,137)}&=\frac65,&
  \frac{s(15,137)}{s(18,137)}&=\frac9{22}.\label{eq:dedekind-five-ratios-b}
\end{align}
The same table also gives
\[
  \frac{s(7,137)}{s(22,137)}=9,
  \quad \frac{s(9,137)}{s(22,137)}=6,
  \quad \frac{s(11,137)}{s(22,137)}=7,
  \quad \frac{s(3,137)}{s(7,137)}=\frac{11}{4}.
\]
\end{proposition}

\begin{proof}
Direct evaluation of the sawtooth definition in rational arithmetic agrees
with the Euclidean recurrence obtained from reciprocity.  Both routes give the
displayed reduced fractions, factorizations, and ratios.
\end{proof}

\begin{proposition}[Symmetry and zeros]
\label{prop:dedekind-symmetry-137}
\label{prop:dedekind-zero-partition-137}
For $1\leq h\leq136$,
\[
  s(137-h,137)=-s(h,137),
  \qquad
  s(h^{-1}\bmod137,137)=s(h,137).
\]
The complete zero set is
\[
  \{h:1\leq h\leq136,\ s(h,137)=0\}=\{37,100\}.
\]
\end{proposition}

\begin{proof}
The two displayed symmetries follow from
Proposition~\ref{prop:dedekind-identities}, since $137$ is prime.  Exact
enumeration of all $136$ residue representatives gives the two zeros.
\end{proof}

\section{Do other primes share the level-137 ratios?}
\label{sec:dedekind-prime-loci}

Fix the argument pairs below and take the target ratios from their values at
$p=137$.  We classify the primes that
satisfy the corresponding arithmetic equations; the final theorem also fixes
an explicit search window for the arguments.  The mathematical mechanism is
the following residue-class quadratic reduction; the selected numerical
targets are examples of that decision procedure, not independently derived
inputs.

\begin{proposition}[Residue-class quadratic reduction]
\label{prop:dedekind-residue-quadratic}
Let $h,p>0$ be coprime and let $r_h$ be the least nonnegative residue of $p$
modulo $h$; when $h>1$, it is a unit residue.
Define
\[
  \beta_h(r_h)=3h+12h\,s(r_h,h),
  \qquad
  Q_{h,r_h}(X)=X^2-\beta_h(r_h)X+h^2+1.
\]
Then
\begin{equation}\label{eq:dedekind-residue-q}
  12hp\,s(h,p)=Q_{h,r_h}(p).
\end{equation}
Consequently, for positive integers $h_1,h_2$, a reduced target $a/b$, and a
prime $p$ coprime to $h_1h_2$ for which $s(h_2,p)\neq0$, the equation
\[
  \frac{s(h_1,p)}{s(h_2,p)}=\frac ab
\]
is equivalent in each unit residue class modulo
$L=\operatorname{lcm}(h_1,h_2)$ to
\begin{equation}\label{eq:dedekind-p-r}
  P_r(p)=b h_2 Q_{h_1,r_1}(p)-a h_1 Q_{h_2,r_2}(p)=0,
\end{equation}
where $r_i\equiv r\pmod{h_i}$.  Each $P_r$ has degree at most two.
\end{proposition}

\begin{proof}
Apply reciprocity \eqref{eq:dedekind-reciprocity} and periodicity to obtain
\eqref{eq:dedekind-residue-q}; clearing the quotient gives
\eqref{eq:dedekind-p-r}.  Primes dividing $h_1h_2$, which lie outside the
reciprocity hypothesis, must instead be evaluated directly from
Definition~\ref{def:classical-dedekind-sum}.
\end{proof}

\begin{theorem}[Five fixed-input prime loci]
\label{thm:dedekind-fixed-pair-prime-loci}
For each row below, the domain is all positive primes $p$ for which the
denominator sum is nonzero.  In every row, the displayed ratio equals its
supplied target if and only if $p=137$.
Put
\[
  Z_2=\{p:s(h_2,p)=0\},\qquad
  Z_{00}=\{p:s(h_1,p)=s(h_2,p)=0\},
\]
and let $Z_{\ne0,0}$ contain the denominator-zero primes having nonzero
  numerator.  The exact classification is
\[
\begin{array}{c|c|c|c|c}
(h_1,h_2,a/b)&\text{defined locus}&Z_2&Z_{00}&Z_{\ne0,0}\\
\hline
(22,7,1/9)  &\{137\}&\{2,5,7\}     &\{2,5\}   &\{7\}\\
(19,9,1/12) &\{137\}&\{2,3,41\}    &\{2\}     &\{3,41\}\\
(13,7,1/4)  &\{137\}&\{2,5,7\}     &\{2,5\}   &\{7\}\\
(14,28,6/5) &\{137\}&\{2,5,7,157\} &\{2,7\}   &\{5,157\}\\
(15,18,9/22)&\{137\}&\{2,3,5,13\}  &\{2,3,5\} &\{13\}
\end{array}
\]
These are complete finite sets of primes, not search cutoffs.
\end{theorem}

\begin{proof}
For the five rows, Proposition~\ref{prop:dedekind-residue-quadratic} produces
respectively $60,108,72,12,$ and $24$ unit-residue equations, $276$ in total.
No equation is the zero polynomial.  Each exact linear or quadratic equation
is solved over $\Q$, then filtered by its residue class, positive integrality,
primality, and the nonzero-denominator condition.  The only retained prime in
each row is $137$.  Every prime dividing an input argument is evaluated
directly, and a separate exact denominator equation gives the complete
undefined sets printed in the table.

For the first row, the following equivalent notation makes the residue
quadratic explicit.  For a
unit residue $r$ modulo $154$, set
\[
  C_r=\frac{63\beta_{22}(r\bmod22)-22\beta_7(r\bmod7)}{1848}.
\]
The residue equation is
\begin{equation}\label{eq:dedekind-ratio-general}
  41p^2-1848C_rp+29455=0.
\end{equation}
For $r=137\bmod154$, the exact values
$\beta_{22}=102$, $\beta_7=27$, and $C_r=243/77$ give
\begin{equation}\label{eq:dedekind-ratio-137}
  41p^2-5832p+29455=0,
\end{equation}
whose rational roots are $137$ and $215/41$.  The factorization
$29455=5\cdot43\cdot137$ concerns the constant term, not the discriminant,
and is a consequence of the integer root rather than independent selection
evidence.  At $p=5$, both $s(22,5)$ and $s(7,5)$ vanish.  Thus $p=5$ solves
only the cross-multiplied equation $0=0$; the quotient is undefined.
\end{proof}

\begin{corollary}[The $(22,7)$ prime locus]
\label{cor:dedekind-ratio-prime-locus}
For every prime $p$ with $s(7,p)\neq0$,
\[
  \frac{s(22,p)}{s(7,p)}=\frac19
  \quad\Longleftrightarrow\quad p=137.
\]
\end{corollary}

\begin{proof}
This is the first row of
Theorem~\ref{thm:dedekind-fixed-pair-prime-loci}, with its quotient domain
stated explicitly.
\end{proof}

\begin{theorem}[An $h\leq28$ target window]
\label{thm:dedekind-h28-target-window}
Let
\[
  \mathcal T=\left\{\frac19,\frac1{12},\frac14,\frac65,\frac9{22}\right\}
\]
and, for a prime $p$, let
\[
  \mathcal R_{28}(p)=
  \left\{
    \frac{s(h_1,p)}{s(h_2,p)}:
    1\leq h_1,h_2\leq28,\ s(h_2,p)\neq0
  \right\}.
\]
Then, for primes $p\geq5$,
\[
  \mathcal T\subseteq\mathcal R_{28}(p)
  \quad\Longleftrightarrow\quad p=137.
\]
\end{theorem}

\begin{proof}
Apply Proposition~\ref{prop:dedekind-residue-quadratic} to the
$5\cdot28\cdot28=3920$ target and ordered-pair choices.  Across the unit
residue classes this gives $252940$ equations: $252712$ quadratic, $223$
linear, and five nonzero constant equations, with no identity equation.
Solving these equations over $\Q$ gives $162$ nonempty choices in their
range of applicability.  Including direct evaluations at the exceptional
small primes $5\leq p\leq\max(h_1,h_2)$ brings the total to $195$
nonempty choices and proves the other $3725$ empty.
The five target-prime sets have sizes $31,12,40,8,$ and $7$;
their intersection is $\{137\}$, and the largest prime in any individual set
is $11813$.

These residue-class equations cover $p>\max(h_1,h_2)$.  For
$5\leq p\leq\max(h_1,h_2)$, including primes dividing an argument, direct
evaluation of the sawtooth definition supplies the exceptional cases.  At
$p=137$, the five pairs in
Theorem~\ref{thm:dedekind-fixed-pair-prime-loci} are explicit witnesses.  The
complete coefficient data and an exact-arithmetic verification routine are
described in Appendix~A.
\end{proof}

\begin{remark}[Interpretation]
Both theorems solve arithmetic equations whose pairs and fractions are
chosen from the level-$137$ data; the second also fixes the cutoff $28$.
Here ``prime locus''
means only the exact solution set under those inputs.
\end{remark}

\section{Conclusion}

The arithmetic of $X_0(137)$ brings several exact structures into a single
picture: genus $11$, Fricke split $(4,7)$, quotient genus $4$, a cyclic
class group of order $8$, and two newform orbits of degrees $4$ and $7$.  The
equality between the Fricke fixed-point count $n_{137}=8$ and the class number
$h(-548)=8$ places the geometric and quadratic-form calculations side by
side.  The rational period map proves rank zero for the degree-seven orbit,
while a rigorous treatment of the degree-four derivative remains an open
next step.

The Hom and Dedekind classifications illustrate a common method: an intrinsic
formula first bounds the possibilities, after which exact finite arithmetic
determines the locus.  It would be useful to study other cross-Fricke targets
uniformly and to formulate broader residue-quadratic theorems for families of
Dedekind ratios.  The present results supply exact test cases for both
directions without claiming that the chosen targets are themselves canonical.

\appendix

\section{Computational data}

Machine-readable supplementary data record the exact $11\times11$ matrix of
$T_2$, the cyclic similarity and irreducibility data, the $23\times14$
rational period map, and its nonzero winding-symbol image.  A second dataset
contains all $3920$ target and ordered-pair cases and all $252940$
unit-residue equations used in Theorem~\ref{thm:dedekind-h28-target-window}.
Verification routines recompute the relevant matrix identities, rational
roots, exceptional primes, and direct Dedekind sums using exact arithmetic.
These files are included in the accompanying package: \texttt{certificates/}
contains the exact data and \texttt{tools/} contains their verifiers.

The smaller modular-symbol and class-group computations were carried out in
SageMath 10.7 with PARI 2.17.3.  All theorem-level matrix, polynomial,
form-class, dimension, Dedekind-value, and residue-root statements use exact
arithmetic.  The high-precision derivative evaluations for the degree-four
orbit are retained only as conditional diagnostics.

\section*{Acknowledgment of assistance}
AI-assisted research, computation, checking, and editorial support used
OpenAI Codex and Anthropic Claude. Luqman Mahmood directed the work
and is responsible for the claims and final text. These internal
checks are not external peer review.

\bibliographystyle{alphaurl}
\bibliography{references}

\end{document}
